% Standalone theorem/proof fragment. Uses amsmath, amssymb, amsthm.
\section{A depth-independent Gaussian result for the simplest targets}
\label{sec:simpletargets}

The norm-capacity obstruction does not automatically make even the simplest
targets hard for the specified Gaussian prior. In fact constants and dictators
have a depth-independent sample bound in the strong-prior limit.

\begin{theorem}
Fix $d\ge1$, $N\ge1$, and $L\ge2$, with tanh hidden activations, the constant
input coordinate, and a linear readout as defined in this report's model. Let
\[
 \mathcal M_d=\{+1,-1,+x_1,-x_1,\ldots,+x_d,-x_d\}.
\]
Among all Boolean functions implemented with the all-input strict margin
$1/2$, the functions of smallest weighted norm are exactly $\mathcal M_d$.
Their common infimum norm value depends on $d,L$ but not on $N$.

Consequently, for every $f\in\mathcal M_d$ and every $\varepsilon,\delta\in(0,1)$,
there is $\sigma_0>0$ such that the original globally margin-conditioned
Gaussian posterior has posterior-draw risk at most $\varepsilon$, with
probability at least $1-\delta$, whenever
\[
 0<\sigma<\sigma_0,\qquad
 n\ge\frac2\varepsilon
       \left(\log(2(d+1))+\log\frac1\delta\right).
\]
The guarantee holds for arbitrary input distributions on the Boolean cube.
The sample bound is independent of depth and learner width.
The variance threshold may depend on $d,L,N,\varepsilon,f$; no uniform
threshold over depths or widths is asserted.
\end{theorem}

\begin{proof}
First, the function $Q(u)=\tanh^2\sqrt u$ is strictly increasing and strictly
concave for $u\ge0$. To see concavity, put $t=\sqrt u>0$ and $y=\tanh t$.
A direct derivative computation gives
\[
 Q''(u)=\frac{1-y^2}{2t^3}\bigl[t(1-3y^2)-y\bigr]<0.
\]
Indeed $y-t(1-y^2)>0$ for $t>0$, since it vanishes at zero and has derivative
$2ty(1-y^2)>0$; hence the bracket is less than $-2ty^2$.
The continuous extension at zero preserves concavity.
Jensen's inequality therefore implies, for every $v\in\mathbb R^N$,
\begin{equation}\label{eq:tanh-rms}
 \frac{\|\tanh v\|_2}{\sqrt N}
 \le\tanh\!\left(\frac{\|v\|_2}{\sqrt N}\right).
\end{equation}

For any network weights define
\[
 a=\frac{\|W^1\|_F}{\sqrt N},\qquad
 b_s=\|W^s\|_F\ (2\le s<L),\qquad
 c=\sqrt N\|v\|_2.
\]
Their weighted squared norm is
\[
 \kappa_N=d a^2+\sum_{s=2}^{L-1}b_s^2+c^2.
\]
Define the scalar recursion $r_1=\tanh a$, $r_s=\tanh(b_s r_{s-1})$ for
$2\le s<L$, and let $F(a,b,c)=c r_{L-1}$.
Let $K_*(d,L)$ be the minimum of the displayed cost over nonnegative
$(a,b_2,\ldots,b_{L-1},c)$ with $F\ge1/2$.
The minimum exists: the constraint is closed and nonempty and the cost is
coercive. Every feasible tuple has all entries positive. At a minimizer,
$F=1/2$, since otherwise decreasing $c$ lowers the cost.

Every tuple is realized by a constant-channel network: all first-layer rows
have weight $a$ on the constant input and zero elsewhere; every entry of
$W^s$ is $b_s/N$; and every readout entry is $c/N$.
Its output is the constant $F$ and its norm is the stated cost, at every width.
Moving the first-layer weight from the constant input to input $x_i$ gives
output $F x_i$ by oddness of tanh; negating the readout gives the negative
functions. Thus all members of $\mathcal M_d$ have infimum norm at most $K_*$,
with strict margins obtained by an arbitrarily small increase of the readout.

Conversely, consider a network with $|z_w(x)|\ge1/2$ at every Boolean input.
For uniform $x$ on the Boolean cube and $\bar x=(1,x)$,
\[
 \mathbb E_x\frac{\|W^1\bar x\|_2^2}{N}
 =\frac{\|W^1\|_F^2}{N}=a^2.
\]
There is therefore an $x_*$ with
$\|W^1\bar x_*\|_2/\sqrt N\le a$.
Apply \eqref{eq:tanh-rms} at each hidden layer, the inequality
$\|Wh\|_2\le\|W\|_F\|h\|_2$ at each internal affine map, and Cauchy--Schwarz
at the readout. These give
\[
 \frac12\le|z_w(x_*)|\le F(a,b,c).
\]
Thus its norm is at least $K_*$. In particular, the members of
$\mathcal M_d$ have exactly the global infimum norm $K_*$.

We check the equality case to exclude all other Boolean functions.
Suppose a globally margin-$1/2$ network has norm exactly $K_*$.
Its associated scalar tuple minimizes the scalar problem, so $F=1/2$.
All parameters in that tuple are positive, so the scalar recursion is strictly
increasing in the first preactivation RMS.
Choose $x_*$ minimizing that RMS. If its value were strictly less than $a$,
then the preceding chain would yield $|z_w(x_*)|<1/2$.
Hence the minimum RMS equals its quadratic mean $a$, so
\[
 \frac{\|W^1\bar x\|_2}{\sqrt N}=a
 \quad\text{for every Boolean }x.
\]
For every input, the entire output bound is consequently an equality:
$|z_w(x)|=F=1/2$. In particular the first hidden vector has the fixed nonzero
norm $\sqrt N\tanh a$.
If $L=2$, equality in readout Cauchy--Schwarz says that all first hidden
vectors are collinear with the same nonzero readout vector.
If $L\ge3$, equality in
$\|W^2h^1(x)\|_2\le\|W^2\|_F\|h^1(x)\|_2$ says the same: every nonzero row
of $W^2$ is collinear with $h^1(x)$, and $W^2\ne0$.
Thus, fixing any reference input $x_0$, there is
$\epsilon(x)\in\{-1,1\}$ with
\[
 h^1(x)=\epsilon(x)h^1(x_0).
\]
Oddness and injectivity of tanh imply
$W^1\bar x=\epsilon(x)W^1\bar x_0$.
Choose a nonzero coordinate of $W^1\bar x_0$; the corresponding first-layer
affine form therefore has constant nonzero absolute value on the cube.
An affine form $\alpha_0+\sum_i\alpha_i x_i$ with constant squared value can
have only one nonzero coefficient: expanding its square forces
$\alpha_0\alpha_i=0$ and $\alpha_i\alpha_j=0$ for $i\ne j$.
Therefore $\epsilon$ is a signed constant or dictator.
Every remaining layer is odd, so $z_w(x)=\epsilon(x)z_w(x_0)$, proving that
its Boolean function belongs to $\mathcal M_d$.

The strict-margin infimum causes no missing equality case.
For each representable Boolean function $g$, replace strict constraints by
$g(x)z_w(x)\ge1/2$. The relaxed feasible set is closed and its minimum norm
is attained by coercivity in the finite-dimensional parameter space.
Its minimum equals the strict infimum: scaling the readout by $1+\eta$ gives
strict feasibility with norm converging to the relaxed value as $\eta\downarrow0$.
Consequently every representable $g\notin\mathcal M_d$ has
$\kappa_N(g)>K_*$.

Finally, the small-variance Gaussian calculation in
Section~\ref{sec:strong-limit}, applied to the open strict-margin
implementation sets $U_{N,g}$, gives
\[
 \frac{\mu_{N,\sigma}(U_{N,g})}
      {\mu_{N,\sigma}(U_{N,f})}\longrightarrow0
 \quad\text{when }\kappa_N(g)>\kappa_N(f).
\]
There are finitely many Boolean functions. Therefore, for $f\in\mathcal M_d$,
\[
 \frac{\pi_{N,\sigma}(\mathcal M_d^c)}{\pi_{N,\sigma}(f)}\longrightarrow0.
\]
Choose $\sigma_0$ so that this ratio is at most $\varepsilon/2$.
At the displayed sample size, a union bound shows that, with probability
$1-\delta$, every consistent member of $\mathcal M_d$ has error at most
$\varepsilon/2$. The posterior mass outside $\mathcal M_d$ is at most
$\pi_{N,\sigma}(\mathcal M_d^c)/\pi_{N,\sigma}(f)\le\varepsilon/2$,
because $f$ is always consistent. Averaging errors proves the claimed
posterior risk guarantee.
\end{proof}

\paragraph{Scope.}
This is a result for the actual independent Gaussian prior in the strong-prior
limit, with the all-input margin conditioning retained. It does not settle
arbitrary circuit-simple targets, the specified nonzero variances in the
$\gamma\ge1$ papers, or a width/depth-uniform variance threshold.
It shows specifically why a uniform-truth-table low-norm prior counterexample
with a constant target cannot simply be transferred to this Gaussian prior.
